\documentclass[11pt]{article}

\usepackage{amsmath}
\usepackage{amssymb}
\usepackage[margin=1in]{geometry}
\usepackage{hyperref}

\title{An information-content ceiling check for the CNP-AIRT ablation ladder}
\author{Working note --- not part of the NeurIPS submission}
\date{\today}

\begin{document}
\maketitle

\section{Purpose}

Every neural variant in the ablation ladder (\texttt{ablation-results.md},
\texttt{uniform-audit.md}) ties or loses to the Beta-Binomial baseline on
calibrated Brier ALC. Before attributing this to undertraining or an
objective bug, check whether the baseline is already near the
\emph{information-theoretic ceiling} for this task: with at most 31 binary
responses, there is a hard limit on how precisely a scalar ability
$\theta$ can be known, and therefore a hard floor on achievable Brier risk.
If the Beta-Binomial baseline already sits at that floor, no architecture
change can beat it, and the ladder rungs are not testing anything.

\section{Model}

Assume a correctly specified 2PL item response model. Item $i$ has
discrimination $a_i > 0$ and difficulty $b_i$:
\begin{equation}
  p_i(\theta) = \sigma\!\big(a_i(\theta - b_i)\big), \qquad
  \sigma(x) = \frac{1}{1+e^{-x}}.
\end{equation}

\section{Step 1: Fisher information from the acquired context}

The Fisher information that a single response on item $i$ carries about
$\theta$ is
\begin{equation}
  I_i(\theta) = a_i^2\, p_i(\theta)\big(1 - p_i(\theta)\big).
\end{equation}
For a context $C$ of $n$ acquired items, the total information at the true
ability is additive:
\begin{equation}
  I_n(\theta) = \sum_{i \in C} a_i^2\, p_i(\theta)\big(1 - p_i(\theta)\big).
\end{equation}
This is exactly the quantity already computed as the negative Hessian in
the existing MAP solver (\texttt{faithful\_np.py}, \texttt{map\_theta}):
\begin{verbatim}
hessian = -sum(a_C^2 * p_C * (1 - p_C)) - 1 / prior_variance
\end{verbatim}
The codebase is already computing $I_n(\theta)$ at every forward pass; it is
simply not being read as a bound.

\section{Step 2: posterior variance (Laplace approximation)}

With a Gaussian prior of variance $\tau^2$ on $\theta$ (the code uses
$\tau^2 = 4$), the Laplace approximation to the posterior at the MAP
estimate $\hat\theta$ gives
\begin{equation}
  \sigma_\theta^2(n) \;\approx\; \frac{1}{I_n(\hat\theta) + 1/\tau^2}.
  \label{eq:posterior-var}
\end{equation}
By the Cram\'er--Rao bound, no estimator --- linear, Bayesian, or a trained
neural network --- can achieve a lower variance than
Eq.~\eqref{eq:posterior-var} if the 2PL model is correctly specified and
$(a_i, b_i)$ are known. This is the hard floor referred to loosely as
``$\sim 31$ bits of information'': the precise statement is that
$\sigma_\theta^2(n)$ is bounded below by the reciprocal of total Fisher
information, not literally Shannon bits.

\section{Step 3: propagating ability uncertainty into item-level Brier risk}

For a target item $t$ with parameters $(a_t, b_t)$, the Bayes-optimal
predictor under squared loss is the posterior-predictive probability
\begin{equation}
  \hat p_t = \mathbb{E}_{\theta \sim q(\theta \mid C)}\big[p_t(\theta)\big].
\end{equation}
Its achievable risk decomposes exactly into an \emph{estimation} term and an
\emph{irreducible noise} term. For any fixed $\theta$,
\begin{equation}
  \mathbb{E}_{y \sim \mathrm{Bernoulli}(p_t(\theta))}\!\left[(\hat p_t - y)^2 \mid \theta\right]
  = (\hat p_t - p_t(\theta))^2 + p_t(\theta)\big(1 - p_t(\theta)\big).
\end{equation}
Taking the expectation over $\theta \sim q(\theta \mid C)$ and using
$\hat p_t = \mathbb{E}_\theta[p_t(\theta)]$ for the first term:
\begin{equation}
  R(t) \;=\; \underbrace{\mathrm{Var}_{\theta \sim q}\big[p_t(\theta)\big]}_{\text{residual ability uncertainty}}
  \;+\; \underbrace{\mathbb{E}_{\theta \sim q}\big[p_t(\theta)\big(1 - p_t(\theta)\big)\big]}_{\text{irreducible Bernoulli noise}}.
  \label{eq:risk-decomp}
\end{equation}
This is the Bayes risk floor for item $t$: it is what the best possible
predictor achieves, given only the information content of the context, with
the true item parameters known exactly.

\subsection{First-order closed form}

Linearizing $p_t(\theta)$ around $\hat\theta$,
$p_t(\theta) \approx \hat p_t + a_t\, \hat p_t(1-\hat p_t)(\theta - \hat\theta)$,
and substituting $\sigma_\theta^2(n)$ from Eq.~\eqref{eq:posterior-var} gives
a closed-form ceiling that needs no simulation:
\begin{equation}
  R(t; n) \;\approx\;
  \hat p_t \big(1 - \hat p_t\big)
  \;+\;
  a_t^2\, \hat p_t^2 \big(1 - \hat p_t\big)^2\, \sigma_\theta^2(n),
  \qquad \hat p_t = p_t(\hat\theta).
  \label{eq:closed-form}
\end{equation}

\section{Step 4: the ceiling to compare against}

Using the competition's own budget grid and weights,
\begin{equation}
  \mathrm{ALC}_{\min} \;=\; \sum_{b \in \{0,1,3,7,15,31\}} w_b \cdot
  \mathbb{E}_{\text{tasks},\, t}\big[R(t;\, n=b)\big],
  \label{eq:alc-floor}
\end{equation}
where the expectation over target items $t$ and tasks is taken either

\begin{itemize}
  \item \textbf{in closed form}, by plugging the corpus's empirical
    distribution of $(a_i, b_i)$ into Eq.~\eqref{eq:closed-form} and
    averaging, or
  \item \textbf{by Monte Carlo}: simulate $\theta_{\mathrm{true}} \sim
    \mathcal{N}(0, \tau^2)$, draw context and target responses from real or
    fitted $(a_i, b_i)$, run the \emph{existing} MAP + Hessian solver
    exactly as in production, evaluate $R(t)$ by quadrature over the
    resulting $q(\theta)$, and average over many simulated tasks.
\end{itemize}

\section{The check}

Compare $\mathrm{ALC}_{\min}$ from Eq.~\eqref{eq:alc-floor} against the
measured Beta-Binomial Brier ALC on validation ($0.222838$, per
\texttt{uniform-audit.md}).

\begin{itemize}
  \item If the measured value is within noise of $\mathrm{ALC}_{\min}$, the
    baseline is already near-optimal given the information content of
    $\le 31$ binary responses. No architecture --- CNP, ANP, GP-residual,
    or latent --- can beat it without violating the Cram\'er--Rao bound,
    and ladder rungs 3 through 6 are testing against an unbeatable number
    rather than against a weak baseline.
  \item If the measured value is meaningfully above $\mathrm{ALC}_{\min}$,
    there is real headroom, and the null results documented in
    \texttt{ablation-results.md} and \texttt{uniform-audit.md} point to a
    training or objective problem (undertraining, KL misweighting, detached
    MAP gradient) rather than to a fundamental information limit.
\end{itemize}

This is a half-day closed-form-plus-Monte-Carlo computation against code
that already exists (\texttt{map\_theta}), and it should be run before
committing further compute to architecture rungs 3--6 of the ablation
ladder.

\end{document}
